\section{Introduction}
\label{sec:introduction}

Large language models are increasingly used for security knowledge work:
explaining weaknesses, drafting detections, reviewing code, summarizing
vulnerabilities, and assisting with incident response. The most capable systems,
however, are commonly served from data centers with enterprise GPUs. Training
and even inference can require hardware that independent researchers, students,
small security teams, and air-gapped organizations do not have. This is not only
an efficiency problem; it shapes who can study, audit, and deploy security AI.

\cyberlite{} investigates a deliberately constrained alternative. The primary
model is Qwen3.5-0.8B, a sub-billion-parameter checkpoint with a hybrid
linear/full-attention architecture and a vision tower that is unused in this
work. The research question is:

\begin{quote}
\emph{Can a 0.8-billion-parameter language model achieve competitive performance
on specialized cybersecurity tasks through parameter-efficient fine-tuning,
knowledge distillation, retrieval augmentation, and quantization, while remaining
practical to train and deploy on consumer hardware?}
\end{quote}

Answering this question requires more than a training script. It requires
licensed and auditable data, contamination control, reproducible configurations,
task-specific metrics, memory and latency measurement, and a clear separation
between hypotheses, engineering validation, and scientific results. This
manuscript and the accompanying repository provide that scaffolding.

\paragraph{Contributions.}
\begin{enumerate}[leftmargin=1.4em]
  \item An architecture-aware, low-VRAM post-training pipeline for Qwen3.5-0.8B
  that explicitly handles its hybrid attention, multimodal checkpoint class,
  LoRA target modules, and fill-in-the-middle tokenizer.
  \item A provenance-first dataset preparation pipeline with license gates,
  secret redaction, exact and near-duplicate removal, deterministic splitting,
  and 13-gram benchmark contamination screening.
  \item A local inference and autocomplete stack with streaming, an
  OpenAI-compatible subset, non-executing completion, and CPU/GPU/MPS paths.
  \item A task-specific evaluation harness with fixed seeds, uncertainty
  intervals, false-positive measurement, lexical grounding proxies, and
  measured memory/latency/cost reporting.
  \item A falsifiable experimental design and a paper scaffold in which every
  performance result is explicitly marked pending.
\end{enumerate}

\paragraph{Status.}
This is a work-in-progress manuscript. The repository contains a working
minimum viable pipeline, but \textbf{no trained research checkpoint and no
verified benchmark result are reported here}. Statements about Qwen3.5's
architecture and parameter count are verified against the checkpoint
configuration and a local load. Statements about future performance are
hypotheses. Numerical tables without measurements are placeholders and must not
be cited as findings.

\paragraph{Reproducibility and safety.}
The default workflow uses no paid APIs, does not download datasets
automatically, and never executes generated commands. The repository records
effective configurations, dataset manifests, environment versions, hardware,
and Git state for every run. These choices are part of the research
contribution, not an afterthought.
